% `tex.web` §882: breaking at a discretionary whose replace text is not empty.
%
% A discretionary is three lists: what goes before a break, what goes after,
% and what stands there when no break is taken. §145 stores the third as
% `replace_count` nodes *following* the disc node, and §882 throws them away
% when the break is taken — the next line begins after the disc node **and**
% after its replace text.
%
% Our line *builder* knew that; it cut at `cut + 1 + d.replace`. Our line
% *searcher* did not — it started the next line with `after_break`, which skips
% only *discardable* nodes, and replace text is usually characters and boxes.
% So the search measured a next line that still had the replace text at its
% head, found the break infeasible, and took no break at all:
%
%     \hsize=20pt   AAAA\discretionary{X}{Y}{Z}BB\par
%     tex    two lines — `AAAA \discretionary X` then `Y … BB`
%     ours   one overfull line holding the whole discretionary
%
% An empty replace text always worked, which is why nothing caught this: the
% two computations only disagree when there is something to throw away, and
% `\-` and every hyphenation break have none.
%
% §882 also leaves the disc node itself in the line, emptied, ahead of its
% pre-break text — that is what a bare `\discretionary` in a box display is.
\catcode`\{=1 \catcode`\}=2
\showboxdepth=3 \showboxbreadth=12 \tracingonline=0
\font\r=cmr10 \r \hsize=20pt \parindent=0pt \rightskip=0pt \tolerance=10000

\setbox0\vbox{AAAA\discretionary{X}{Y}{Z}BB\par}
\showbox0

% several replace nodes, all thrown away
\setbox1\vbox{AAAA\discretionary{X}{Y}{ZZZ}BB\par}
\showbox1

% a box counts as one node however wide
\setbox2\vbox{AAAA\discretionary{X}{Y}{\hbox to 5pt{}}BB\par}
\showbox2

% an empty replace text — the case that always worked, kept so it stays working
\setbox3\vbox{AAAA\discretionary{X}{Y}{}BB\par}
\showbox3

% and one that does not break, where the replace text is what gets set
\hsize=200pt
\setbox4\vbox{AA\discretionary{X}{Y}{Z}BB\par}
\showbox4
\end
